\chapter{How to use the templates} \label{ch-template}

This is a practical guide into how to use this template, by explaining the role of the different folders, and an option of \verb|\documentclass{oist_thesis}|, which accepts either \verb|temporary| or \verb|final|.

\section{Folders}

The main folder contains three folders detailed here:

\begin{itemize}

	\item \textbf{Images.} This folder should contain all the images that you will use in your thesis. It can contain subfolders, for example one for each chapter. To include an image from the main text, use something like \texttt{\textbackslash includegraphics\{subfolder/image.jpg\} } without worrying about the \texttt{Images} path.

	\item \textbf{MainText.} This folder contains a series of \LaTeX\ files that form the main text: introduction, chapters, conclusion, appendices and published articles. The introduction and conclusion as they are now are not numbered, which creates a few difficulties with the headers of the thesis. Those are solved by including the commands \texttt{\textbackslash unnumberedchapter\{\}} and \texttt{\textbackslash numberedchapter} before including the files in \texttt{Thesis.tex}. If you want the introduction and conclusion to be numbered, include them as regular chapters.

	\item \textbf{Preamble.} This folder contains a series of \LaTeX\ files with the pages that will appear before the main text. Please write (or copy and paste) your own text in those files and delete the dummy text when appropriate. The files are:
	      \begin{itemize}
		      \item \texttt{abbreviations.tex} --- List of abbreviations. If the list goes over one page, create another table.
		      \item \texttt{abstract.tex} --- Abstract. Follow directions in the file.
		      \item \texttt{acknowledgments.tex} --- Acknowledgments. Follow directions in the file.
		      \item \texttt{declaration.tex} --- Declaration of Original and Sole Authorship. Only modify the last item. This page needs to be signed once printed.
		      \item \texttt{dedication.tex} --- Dedication (optional). Should only be a very few lines.
		      \item \texttt{glossary.tex} --- Glossary (optional). If the list goes over one page, create another table.
		      \item \texttt{nomenclature.tex} --- Nomenclature (optional). If the list goes over one page, create another table.
		      \item \texttt{physics\_bibstyle.bst} --- Bibliography style file modified by Jeremie Gillet in 2011 to suit his thesis. Might be suitable for physics. If you want to use another custom bibliography style, include the file in this folder.
		      \item \texttt{Thesis\_bibliography.bib} --- BibTeX file containing your bibliography.
		      \item \texttt{Thesis\_publications.bib} --- BibTeX file containing the list of your own publications (only for thesis by publication).
	      \end{itemize}

\end{itemize}

\section{\texttt{Thesis.tex}}

This is the main file, the only one that needs to be compiled to build the thesis. Compile once with \LaTeX, once with BibTeX and finally twice with \LaTeX\ to get all the references right. At the top of this file, you can see \verb|\documentclass[temporary]{oist_thesis}|. When you submit a temporary version to the graduate school, do not modify it. When you submit a final version, use \verb|\documentclass[final]{oist_thesis}| instead.

Let's go through each section and comment them briefly. The last section will emphasize the differences between options \verb|\documentclass[temporary]{oist_thesis}| and \verb|\documentclass[final]{oist_thesis}|.

\subsection{REQUIRED PACKAGES AND CONFIGURATIONS}

This section contains the minimum number of packages and definitions to compile the thesis. No line should be removed or modified.

\subsection{ADD YOUR PACKAGES}

Add you packages in this section.
In general, try to include as little packages as possible, as some of them can interfere and break each other.

\subsection{ADD YOUR DEFINITIONS AND COMMANDS}

Optionally, add custom definitions in this section.
You don't have to keep the examples that come with the template.

\subsection{PICK YOUR BIBLIOGRAPHY STYLE}

Pick the bibliography style the most adapted to your field within the \texttt{natbib} package.

\subsection{FOR THESIS BY PUBLICATION ONLY}

There are two sections by that name, which are meant for students publishing a thesis by publication only.
If you are not publishing a thesis by publication, please remove the sections.

The \texttt{multibib} package allows more than one list of references, but extra ones need to be given a name, the \texttt{Pub} in \verb|\newcites{Pub}{List of Publications}|. With this, you may use new commands like \texttt{\textbackslash bibliographystylePub} or \texttt{\textbackslash citePub}. Please update the list of your publications in \texttt{Preamble/Thesis\_publications.bib}.

\subsection{TITLE PAGE}

Creates the title page. Do not modify.

\subsection{PREAMBLE PAGES}

Structures the style (header) for the preamble pages and builds them. Do not modify, except for deleting the optional preambles you might not want to include.

\subsection{LIST OF CONTENTS/FIGURES/TABLES}

Creates the different lists. Do not modify.

\subsection{THESIS MAIN TEXT}

Structures the style for the main text chapters and builds them.

The command \texttt{\textbackslash numberedchapter} is only relevant for a transition between unnumbered sections and numbered sections, it does not need to be included between each chapter.

\subsection{BIBLIOGRAPHY}

Includes the bibliography.

\subsection{APPENDICES}

Structures the style for the appendices and builds them. The appendices are numbered with letters but are structured like regular chapters.

\subsection{Differences between temporary version and final version}

There is only one difference between \verb|\documentclass[temporary]{oist_thesis}| and \verb|\documentclass[final]{oist_thesis}|: the line spacing is different.
This might create small issues, such as page breaking with large tables, images or captions, when compiling the same content.
